\section{学术综述与研究空白}

\begin{frame}{综述协议：范围综述，不伪装成 PRISMA 系统综述}
\scriptsize
\begin{columns}[T,onlytextwidth]
\column{0.48\textwidth}
\begin{block}{检索与纳入}
\begin{itemize}
  \item 截止：2026-07-15
  \item arXiv/作者页、OpenReview、IEEE/DOI
  \item NIST、OWASP、MITRE、OASIS 与源官方文档
  \item 纳入原始研究、正式标准、官方数据文档
  \item 预印本显式标注；排除产品宣传与不可核验二手材料
\end{itemize}
\end{block}
\column{0.48\textwidth}
\begin{block}{四条研究流}
\begin{enumerate}
  \item agent reasoning/tool use/evaluation
  \item active search/bandit/closed-loop control
  \item LLM/agent security taxonomy and benchmarks
  \item CTI collection, standards and knowledge graphs
\end{enumerate}
最终保留 20 个核心来源。
\end{block}
\end{columns}
\smallsource{完整检索词、矩阵与注释书目：\pathcode{report/literature_review.md}}
\end{frame}

\begin{frame}{综述证据边界：新兴 agent 研究与成熟基础理论分开解读}
\small
\begin{columns}[T,onlytextwidth]
\column{0.48\textwidth}
\begin{block}{时间结构}
2024--2026 工作占比较高，反映 agentic search 与 agent security 的快速演进；同时保留 2002/2003/2012 的 UCB、MAPE-K 与 active search 基础工作，避免把“新”误当“已验证”。
\end{block}
\column{0.48\textwidth}
\begin{block}{证据等级}
正式论文/标准用于方法与定义；预印本用于识别前沿方向；官方 API 文档只证明源能力，不证明算法效果。所有性能主张最终都要由本项目冻结 benchmark 与人工 gold 重测。
\end{block}
\end{columns}

\begin{alertblock}{范围限制}
本综述服务于系统定位和实验设计，不声称完成数据库穷尽检索、双人筛选或发表偏倚评估；因此称为范围综述，而非 PRISMA 系统综述。
\end{alertblock}
\end{frame}

\begin{frame}{研究版图：四条成熟链条尚未在一个系统中闭环}
\scriptsize
\begin{tabularx}{\textwidth}{p{0.17\textwidth} p{0.23\textwidth} X X}
\toprule
研究流 & 代表工作 & 已解决 & 对本项目留下的空白 \\
\midrule
Agent & ReAct、AgentBench、AgentDojo & 推理-行动闭环、多环境能力、utility/security 同测 & 异质 CTI 来源、覆盖合同、可恢复 campaign \\
Search & Search-R1、active search、UCB & 多轮 query、固定预算发现、探索/利用 & source applicability、非平稳 reward、证据独立性 \\
Security & OWASP、MITRE ATLAS、NIST RMF & 风险分类、攻击技术、治理/TEVV & 可计算 topic-channel target 与运行级 stop risk \\
CTI/KG & STIX、CTINexus、AttacKG+、KnowCTI、ThreatKG & 标准表达、抽取、实体对齐、攻击图 & 上游主动采集与 KG gap 反向查询 \\
\bottomrule
\end{tabularx}

\vspace{0.7em}
\begin{center}
\textcolor{Purple}{论文机会 = 把四条链的“接口缺口”形式化并实验验证，而不是再做一个通用 agent 框架。}
\end{center}
\end{frame}

\begin{frame}{Agentic search 文献：重点从“会调用工具”转向“如何评价过程”}
\small
\begin{columns}[T,onlytextwidth]
\column{0.49\textwidth}
\begin{block}{方法脉络}
\begin{itemize}
  \item ReAct：reasoning/action 交错，环境反馈修正计划 \cite{yao2023react}
  \item Search-R1：用结果 reward 学习多轮 search，而非只在推理时提示 \cite{jin2025searchr1}
  \item Active Search：固定预算最大化发现目标类的数量 \cite{garnett2012active}
  \item UCB：在未知收益 arm 间平衡探索与利用 \cite{auer2002ucb}
\end{itemize}
\end{block}
\column{0.47\textwidth}
\begin{block}{评测脉络}
\begin{itemize}
  \item AgentBench：8 个交互环境，暴露长期推理/决策问题 \cite{liu2024agentbench}
  \item DeepResearch Bench：100 任务、22 学科；评价有效引用和引用准确率 \cite{du2025deepresearchbench}
  \item AgentDojo：97 任务、629 安全用例，同测 utility 与 prompt injection \cite{debenedetti2024agentdojo}
\end{itemize}
\end{block}
\end{columns}

\begin{alertblock}{转化到本项目}
最终 item 质量、过程 trace、来源引用、策略安全和预算效率必须一起评价；单一 final answer 或 item count 不足以说明 agent 可用。
\end{alertblock}
\end{frame}

\begin{frame}{CTI 与安全文献：结构化能力成熟，上游获取仍是开放问题}
\small
\begin{columns}[T,onlytextwidth]
\column{0.48\textwidth}
\begin{block}{风险与标准}
\begin{itemize}
  \item OWASP 2025：应用级十大 LLM 风险 \cite{owaspllm2025}
  \item MITRE ATLAS：AI tactic/technique/mitigation/case studies \cite{mitreatlas}
  \item NIST AI RMF/GAI Profile：Govern/Map/Measure/Manage 与 TEVV \cite{nistairmf2023,nistgai2024}
  \item STIX 2.1：CTI 对象关系与交换 \cite{stix21}
\end{itemize}
\end{block}
\column{0.48\textwidth}
\begin{block}{知识图谱}
\begin{itemize}
  \item CTINexus：ICL、层次实体对齐、长关系补全，150 报告/10 平台 \cite{cheng2025ctinexus}
  \item AttacKG+：时间展开攻击图 + TTP \cite{zhang2024attackgkplus}
  \item KnowCTI：领域知识增强实体/关系抽取 \cite{knowcti2024}
  \item ThreatKG：多源 OSCTI 持续采集与入图 \cite{threatkg2022}
\end{itemize}
\end{block}
\end{columns}

\begin{alertblock}{共同缺口}
这些工作多数优化“文档到图”的质量；很少用图的结构空洞、不确定边和 provenance 冲突反向决定下一轮去哪个源、查什么、何时停止。
\end{alertblock}
\end{frame}

\begin{frame}{研究空白 -> 可检验贡献}
\scriptsize
\begin{tabularx}{\textwidth}{p{0.26\textwidth} X X}
\toprule
空白 & 本项目可贡献 & 可验证证据 \\
\midrule
结构化覆盖缺口 & topic-channel-quality evidence tensor & 人工标注 benchmark 上 coverage/precision \\
跨源 reward 不可比 & source-aware contextual/constraint bandit & 同预算 A/B + regret/reward calibration \\
停止无风险语义 & confidence/risk-bounded stopping & premature-stop 与 waste 分解 \\
测量误差被忽略 & relevance/topic/event MSA + uncertainty propagation & 双标注、Kappa、ECE、bootstrap CI \\
采集与 KG 单向 & KG-gap-driven query generation & 图完整度、冲突率、下游问答/调查任务 \\
缺少公开过程数据 & trace-level LLM-security CTI benchmark & 任务、快照、mock API、seed、成本清单 \\
\bottomrule
\end{tabularx}

\end{frame}

\begin{frame}{研究定位：把四条链的接口缺口变成一篇可证伪论文}
\begin{center}
\Large\textcolor{Purple}{\textbf{EvidenceGap-Agent}}\\[0.4em]
\large Risk-bounded Active Intelligence Collection for LLM Security
\end{center}

\vspace{0.8em}
\begin{columns}[T,onlytextwidth]
\column{0.48\textwidth}
\begin{block}{论文核心}
用质量加权的有效证据状态表示“还缺什么”，以 source-aware 策略主动查询，并在覆盖、精度、风险与预算约束下停止。
\end{block}
\column{0.48\textwidth}
\begin{block}{区别于通用 agent}
贡献不依赖更大的模型，而来自问题形式化、证据张量、约束策略、停止条件、过程 benchmark 与可复现实证。
\end{block}
\end{columns}

\begin{alertblock}{最重要的审稿问题}
新指标是否真的对应分析师价值？新策略是否在同预算、同源快照、强基线和人工 gold 上显著优于当前实现？
\end{alertblock}
\end{frame}
